\section{Related Work}

\subsection{Survival Analysis with Clinical Features}
Cox proportional hazards~\cite{coxph_original} remains the dominant survival
model in clinical prediction due to its interpretability and calibration
properties.
Applied to structured EHR features, CoxPH baselines on large biobank
cohorts typically achieve Harrell C-index in the 0.60--0.68 range for
all-cause mortality, with performance driven by biomarkers such as
HbA1c, cystatin~C, and CRP~\cite{harrellcindex}.
Deep survival models (DeepHit~\cite{deephit}, DeepSurv) offer flexibility
but rarely outperform well-regularised CoxPH on structured tabular data
without imaging or genomic inputs.

\subsection{LLM-Based Clinical Prediction}
Foundation models pre-trained on biomedical text have been fine-tuned for
clinical prediction tasks including readmission, sepsis onset, and length
of stay.
Delphi~\cite{delphi2024} pioneered the autoregressive EHR paradigm:
ICD codes and demographic tokens are serialised and processed by a causal
LLM trained to predict future diagnosis tokens, with the death-token
logit repurposed as a risk score.
Concurrent work on encoder-style LLMs (BioClinicalBERT, GatorTron,
Med-PaLM) has shown that dense embeddings of clinical notes and discharge
summaries carry prognostic signal when combined with shallow survival
heads.
However, most evaluations are limited to in-hospital outcomes on small
cohorts and do not systematically disentangle which structural modalities
--- temporal ordering, numeric precision, or diagnostic co-occurrence ---
drive embedding quality.

\subsection{Structured Multimodal Clinical Representations}
The PRCV topic of structured multimodal intelligence motivates models
that go beyond statistical semantics to capture spatial, temporal, and
causal structure in scientific and clinical data.
In the EHR domain, temporal structure (sequence of diagnoses, drug
exposures, lab trends) has been encoded via recurrent networks, temporal
convolutions, and more recently transformer-based architectures trained
on ICD event streams.
Biomarker-level structure --- where numeric values carry clinical meaning
relative to reference ranges and comorbidity context --- has received less
attention in the LLM embedding literature.
Our work directly tests whether a general-purpose LLM encoder can recover
both types of structure from prose representations, and compares this
against explicit temporal sequence encoding.
